\subsubsection{Simulated rat spike data in a navigation task}
\label{subsubsection:synthetic_results_details}

The virtual rat moved along a 1-m linear track, with its velocity sampled from an Ornstein--Uhlenbeck (OU) process:
\begin{equation}
dv_t = \theta(\mu-v_t)\,dt + \sigma\,dW_t,
\end{equation}
with parameters $\theta=1.0$ (mean reversion), $\mu=0.0$ m/s (long-run mean), $\sigma=0.4$ (volatility), and $v_0=0.0$ (initial velocity). The simulation lasted 10 minutes (600 s) at a temporal resolution of 100 Hz.

Place fields were modeled as Gaussian functions of position. Two cells had broad fields ($\sigma=20$ cm) centered at 20 and 80 cm, while two had narrow fields ($\sigma=10$ cm) centered at 30 and 90 cm. Peak firing rates were 20, 16, 18, and 15 Hz, respectively, with a 0.1-Hz baseline. The 96 noise neurons had stable mean firing rates sampled uniformly between 1 and 5 Hz and were independent of position. Spikes were generated as Bernoulli draws using the firing-rate-by-time-step probability.

The resulting spike-count matrix was used to train an SED with 128 latents and a sparsity constraint of four active latents per time bin.
\subsubsection{Real macaque spike data in a reaching task}
\label{subsubsection:churchland_results_additional}

\paragraph{Dataset info}
\begin{itemize}

    \item Exp metadata \& info
    
    \begin{itemize}
        \item Session duration
        \item Stimulus conditions
    \end{itemize}
    
    \item Neural data info
    
    \begin{itemize}
        \item Brain regions
        \item Number of units
        \item Number of spikes
        \item Spike summary stats
    \end{itemize}

    \item Access instructions

\end{itemize}

\paragraph{Latent comparisons with sparseNMF, t-SNE, and UMAP}

\subsubsection{Real mouse spike data in an ethological foraging task}
\label{subsubsection:aeon_results_additional}

\paragraph{Dataset info}
\begin{itemize}

    \item Exp metadata \& info
    
    \begin{itemize}
        \item Session duration
        \item Stimulus conditions
    \end{itemize}
    
    \item Neural data info
    
    \begin{itemize}
        \item Brain regions
        \item Number of units
        \item Number of spikes
        \item Spike summary stats
    \end{itemize}

    \item Access instructions

\end{itemize}

\paragraph{Latent comparisons with sparseNMF, t-SNE, and UMAP}
